% 3.2 Money without banks
\subsection{Money without banks}
\label{sec:money}

If the bottle light is the parable, M-Pesa is the case study. Launched in Kenya
in 2007 by the mobile operator Safaricom, it let people store money on their
phone and send it to anyone else by text message, cashing in and out through
small shops that acted as agents. It began as something narrower: a pilot to
let microfinance borrowers repay loans by phone, funded in part by a
challenge fund at the UK's Department for International Development. Once it
was in people's hands, they used it for something else entirely, sending money
to relatives \parencite{hughes2007}.

Here the removed default is the bank branch. In rich countries, keeping and
moving money runs through a dense network of branches, cheques and cards. Kenya
in the mid-2000s had a thin version of that network, mostly in towns. What it
did have, increasingly, was mobile phones. The reframing was to stop asking
how to bring branches to people and to ask instead what people needed from a
bank: somewhere safe to hold a little money and a way to send it. A phone and a
corner shop could do both.

The scale of what followed is easy to see in survey data
(\cref{fig:findex-accounts}). By \findexYear{}, \kenMobileMoney{} of Kenyan
adults had a mobile money account. The share with an account at a bank or
other financial institution was \kenBankAccount{}, and had actually
\emph{fallen} from \kenBankAccountPeak{} in \kenBankAccountPeakYear{}.
Mobile money did not lead people into the banking system so much as provide a
substitute for it. Across sub-Saharan Africa the change is slower but moving in
the same direction: mobile money accounts went from \ssaMobileMoneyStart{} of
adults in \ssaMobileMoneyStartYear{} to \ssaMobileMoney{} in \findexYear{},
overtaking financial-institution accounts (\ssaBankAccount{}) for the first
time. Globally the industry now counts \gsmaRegisteredAccounts{} registered
accounts, of which \gsmaActiveAccounts{} were used in the last 30 days, moving
more than \gsmaTransactionValue{} in a year \parencite{gsma2026sotir}. That is
a different kind of more.

\begin{figure}[htbp]
\centering
% Figure: account ownership by type, Kenya and sub-Saharan Africa.
% Data: data/processed/findex_accounts.csv (scripts/process_findex.py).
% Findex surveys are years apart, so points are marked and joined only as a guide.
\begin{tikzpicture}
\begin{groupplot}[
  group style={group size=2 by 1, horizontal sep=0.15\linewidth},
  scarcity axis, year axis,
  width=0.47\linewidth, height=0.42\linewidth,
  xmin=2011, xmax=2024, xtick={2011,2014,2017,2021,2024},
  ymin=0, ymax=100, ytick={0,25,50,75,100},
  unbounded coords=jump,
  title style={font=\sffamily\footnotesize, anchor=south west, at={(0,1.02)}, align=left},
]
  \nextgroupplot[title={Kenya (\% of adults)}]
  \addplot[series context] table[x=year, y=ken_any, col sep=comma] {../data/processed/findex_accounts.csv}
    node[end label, pos=1, anchor=south west, text=stone] {Any};
  \addplot[series primary, mark=*, mark size=1.6pt] table[x=year, y=ken_mm, col sep=comma] {../data/processed/findex_accounts.csv}
    node[end label, pos=1, anchor=north west, text=indigo] {Mobile\\money};
  \addplot[series accent, mark=square*, mark size=1.4pt, mark options={solid}] table[x=year, y=ken_fi, col sep=comma] {../data/processed/findex_accounts.csv}
    node[end label, pos=1, text=amber!70!ink] {Bank};

  \nextgroupplot[title={Sub-Saharan Africa (\% of adults)}]
  \addplot[series context] table[x=year, y=ssa_any, col sep=comma] {../data/processed/findex_accounts.csv}
    node[end label, pos=1, text=stone] {Any};
  \addplot[series primary, mark=*, mark size=1.6pt] table[x=year, y=ssa_mm, col sep=comma] {../data/processed/findex_accounts.csv}
    node[end label, pos=1, anchor=south west, text=indigo] {Mobile\\money};
  \addplot[series accent, mark=square*, mark size=1.4pt, mark options={solid}] table[x=year, y=ssa_fi, col sep=comma] {../data/processed/findex_accounts.csv}
    node[end label, pos=1, anchor=north west, text=amber!70!ink] {Bank};
\end{groupplot}
\end{tikzpicture}

\caption{\takeaway{In Kenya, mobile money reaches nearly twice as many adults as
banks do; across the region, it has just overtaken them.} In Kenya,
\kenMobileMoney{} of adults had a mobile money account in \findexYear{}, while
bank or other financial-institution accounts fell from \kenBankAccountPeak{} in
\kenBankAccountPeakYear{} to \kenBankAccount{}. Across sub-Saharan Africa,
mobile money accounts (\ssaMobileMoney{}) overtook financial-institution
accounts (\ssaBankAccount{}) in \findexYear{}, up from \ssaMobileMoneyStart{} in
\ssaMobileMoneyStartYear{}. Mobile money was not asked about in 2011.
Categories overlap: many adults hold both.
\figsource{World Bank, Global Findex Database 2025, series account.t.d,
fiaccount.t.d and mobileaccount.t.d; sub-Saharan Africa excludes high-income
economies \parencite{worldbank_findex}.}}
\label{fig:findex-accounts}
\end{figure}

Did it make people better off? The best-known study says yes.
\textcite{suri2016} followed Kenyan households from 2008 to 2014 and compared
those who ended up living near more M-Pesa agents with those who did not. They
estimated that access to mobile money lifted \mpesaHouseholdsLifted{}
households, about \mpesaHouseholdsShare{} of all Kenyan households, out of
poverty, with the largest effects in households headed by women, and that
around \mpesaWomenMoved{} women moved from farming into business or retail
work. The mechanism, they argue, is resilience: when a household is hit by
illness or a bad harvest, relatives elsewhere can send help the same day.

Three conditions made this possible, and none of them was scarcity itself.
Enough people had phones: in 2006, \kenyaPhoneOwners{} of Kenyans owned a
mobile phone and as many again could use someone else's \parencite{afi2010}.
The regulator left room: the Central Bank of Kenya concluded that it had no
clear authority over a non-bank money-transfer service, and rather than
blocking it, issued a letter of no objection in February 2007, a month before
launch, on condition that Safaricom report to it regularly
\parencite{afi2010}. And there was capital and a platform: Safaricom was part
of the Vodafone group, and the first experiment was paid for jointly by
Vodafone and a donor \parencite{hughes2007,phillips2013dfid}. Today
sub-Saharan Africa accounts for \gsmaRegisteredSSA{} of the world's
registered mobile money accounts and \gsmaValueSSA{} of the value moved
through them \parencite{gsma2026sotir}.


\begin{counterpoint}
The poverty estimate is contested. \textcite{bateman2019} argue that the
\textcite{suri2016} analysis has methodological problems and ignores what
happened to the businesses that lost out, to debt taken on through
phone-based lending, and to the profits that flowed to Safaricom and its
shareholders rather than to users. I have not found a published reply from the
original authors. I do not think the critique overturns the basic fact that
mobile money spread to most Kenyan adults and made sending money cheaper and
safer. But it does mean the precise poverty figure should be treated as an
estimate under debate, not a settled result. And it raises a question I take
up in \cref{sec:against}: how much of this is a story about scarcity at all,
and how much is a story about a well-capitalised company, a patient regulator
and a donor willing to fund the first experiment?
\end{counterpoint}
